\documentclass[a4paper,11pt]{article}

% Packages pour l'encodage, la langue et les images
\usepackage[utf8]{inputenc}
\usepackage[T1]{fontenc}
\usepackage[french]{babel}
\usepackage{graphicx}
\usepackage{float}

% Configuration des marges
\usepackage{geometry}
\geometry{hmargin=2.5cm, vmargin=2.5cm}

% Packages pour le code et les liens
\usepackage{listings}
\usepackage{xcolor}
\usepackage{hyperref}
\usepackage{array}
\usepackage{tabularx}
\usepackage{booktabs}

% Configuration des liens
\hypersetup{
    hidelinks,
    pdfborder={0 0 0},
    linktoc=all,
    pdftitle={TP GLPI - Gestion d'incidents},
    pdfauthor={MIGNARD Maël}
}

% Configuration des blocs de code
\definecolor{codegreen}{rgb}{0,0.6,0}
\definecolor{codegray}{rgb}{0.5,0.5,0.5}
\definecolor{backcolour}{rgb}{0.95,0.95,0.92}

\lstset{
    backgroundcolor=\color{backcolour},
    commentstyle=\color{codegreen},
    basicstyle=\ttfamily\footnotesize,
    breaklines=true,
    frame=single,
    numbers=left,
    numberstyle=\tiny\color{codegray},
    keepspaces=true
}

\renewcommand{\arraystretch}{1.3}

% Informations du rapport
\title{\textbf{TP -- Mise en place d'un mini helpdesk avec GLPI}\\Gestion des incidents et ticketing}

\author{MIGNARD Maël \\ BTS SIO -- Support informatique}

\date{\today}

\begin{document}

\maketitle
\tableofcontents
\newpage

% ============================================================
\section{Introduction}
% ============================================================

L'objectif de ce TP est de mettre en place un mini helpdesk à partir de GLPI sur un environnement Debian avec la pile LAMP (Linux, Apache, MariaDB, PHP). Il s'agit de simuler le fonctionnement d'un service support informatique : installation de l'outil, création d'utilisateurs avec des rôles distincts, traitement de tickets d'incidents et de demandes, puis capitalisation des solutions dans une base de connaissances.

Ce compte rendu présente chaque étape réalisée, accompagnée de captures d'écran et des réponses aux questions posées.

\newpage

% ============================================================
\section{Exercice 1 -- Préparation de l'environnement de travail}
% ============================================================

\subsection{Travail réalisé}

Avant d'installer GLPI, j'ai vérifié que la pile LAMP était opérationnelle sur Debian. Les services Apache et MariaDB ont été installés et activés :

\begin{lstlisting}[language=bash]
sudo apt install apache2 mariadb-server php
sudo systemctl enable --now apache2
sudo systemctl enable --now mariadb
\end{lstlisting}

Pour confirmer que le serveur web et PHP fonctionnent correctement, j'ai accédé à l'URL de GLPI dans le navigateur. La page d'accueil de l'assistant d'installation s'est affichée, confirmant qu'Apache sert bien les fichiers PHP :

\begin{figure}[H]
    \centering
    \includegraphics[width=0.85\textwidth, keepaspectratio]{screen/1.png}
    \caption{Page d'accueil GLPI dans le navigateur -- Apache et PHP opérationnels}
\end{figure}

\subsection{Versions relevées}

\begin{center}
\begin{tabular}{|l|c|}
\hline
\textbf{Composant} & \textbf{Version} \\
\hline
Apache  & 2.4.x \\
PHP     & 8.x.x \\
MariaDB & 10.x.x \\
\hline
\end{tabular}
\end{center}

\textit{(À compléter avec les versions réelles : \texttt{apache2 -v}, \texttt{php -v}, \texttt{mariadb --version}.)}

\subsection{Rôle de chaque composant}

\begin{description}
    \item[Apache :] Serveur web qui reçoit les requêtes HTTP du navigateur et sert les fichiers PHP de l'application GLPI.
    \item[PHP :] Langage de programmation côté serveur utilisé par GLPI pour générer les pages dynamiquement.
    \item[MariaDB :] Système de gestion de base de données qui stocke toutes les informations de GLPI (utilisateurs, tickets, matériel\ldots).
    \item[Navigateur web :] Interface cliente permettant d'accéder à GLPI via une adresse HTTP locale.
\end{description}

\subsection{Réponses aux questions}

\begin{description}
    \item[Q1 -- Pourquoi GLPI a-t-il besoin d'un serveur web ?]
    GLPI est une application web développée en PHP. Pour être exécutée, elle nécessite un serveur web (Apache) qui interprète le code PHP et retourne les pages au navigateur de l'utilisateur. Sans serveur web, les fichiers PHP ne sont que du texte brut.

    \item[Q2 -- Quel est le rôle de la base de données dans un outil de ticketing ?]
    La base de données stocke de façon permanente toutes les informations de GLPI : les comptes utilisateurs, les tickets créés, leur état, les historiques d'interventions, les équipements inventoriés\ldots Sans base de données, aucune information ne serait conservée entre deux sessions.

    \item[Q3 -- Pourquoi utilise-t-on un navigateur pour accéder à GLPI ?]
    GLPI est une application web : son interface est générée en HTML/CSS/JavaScript par le serveur. Le navigateur est le client universel qui permet d'afficher et d'interagir avec cette interface, sans avoir à installer de logiciel supplémentaire sur les postes des utilisateurs.
\end{description}

\newpage

% ============================================================
\section{Exercice 2 -- Installation de GLPI}
% ============================================================

\subsection{Installation des dépendances PHP}

GLPI nécessite plusieurs extensions PHP supplémentaires :

\begin{lstlisting}[language=bash]
sudo apt install php-{apcu,cli,common,curl,gd,mysql,xml,mbstring,bcmath,intl,zip,bz2}
\end{lstlisting}

\subsection{Déploiement des fichiers GLPI}

J'ai téléchargé l'archive GLPI et extrait les fichiers dans le répertoire web :

\begin{lstlisting}[language=bash]
wget https://github.com/glpi-project/glpi/releases/download/10.0.x/glpi-10.0.x.tgz
sudo tar -xzvf glpi-10.0.x.tgz -C /var/www/html/
sudo chown -R www-data:www-data /var/www/html/glpi/
\end{lstlisting}

\subsection{Création de la base de données}

Après connexion à MariaDB, j'ai créé la base et l'utilisateur dédié :

\begin{lstlisting}[language=SQL]
CREATE DATABASE glpi CHARACTER SET utf8mb4 COLLATE utf8mb4_unicode_ci;
CREATE USER 'glpiuser'@'localhost' IDENTIFIED BY 'MotDePasseSecure123!';
GRANT ALL PRIVILEGES ON glpi.* TO 'glpiuser'@'localhost';
FLUSH PRIVILEGES;
\end{lstlisting}

\subsection{Informations de connexion à la base}

\begin{center}
\begin{tabular}{|l|l|}
\hline
\textbf{Paramètre} & \textbf{Valeur} \\
\hline
Serveur & localhost \\
Base de données & glpi \\
Utilisateur & glpiuser \\
Mot de passe & (non divulgué) \\
\hline
\end{tabular}
\end{center}

\subsection{Assistant d'installation web}

J'ai lancé l'assistant depuis le navigateur en accédant à \texttt{http://localhost/glpi/public}, puis j'ai suivi les étapes suivantes.

\subsubsection*{Sélection de la langue}

\begin{figure}[H]
    \centering
    \includegraphics[width=0.85\textwidth, keepaspectratio]{screen/2.png}
    \caption{Sélection de la langue -- Français}
\end{figure}

\subsubsection*{Acceptation de la licence}

\begin{figure}[H]
    \centering
    \includegraphics[width=0.85\textwidth, keepaspectratio]{screen/3.png}
    \caption{Licence GNU GPL v3}
\end{figure}

\subsubsection*{Choix du type d'installation}

\begin{figure}[H]
    \centering
    \includegraphics[width=0.85\textwidth, keepaspectratio]{screen/4.png}
    \caption{Choix \og{}Installer\fg{} pour une nouvelle installation}
\end{figure}

\subsubsection*{Étape 0 -- Vérification de la compatibilité de l'environnement}

L'assistant vérifie que toutes les extensions PHP requises sont présentes.

\begin{figure}[H]
    \centering
    \includegraphics[width=0.85\textwidth, keepaspectratio]{screen/5.1.png}
    \caption{Étape 0 -- Vérification des prérequis (partie 1)}
\end{figure}

\begin{figure}[H]
    \centering
    \includegraphics[width=0.85\textwidth, keepaspectratio]{screen/5.2.png}
    \caption{Étape 0 -- Vérification des prérequis (partie 2) -- tous les tests passent}
\end{figure}

\subsubsection*{Étape 1 -- Configuration de la connexion à la base de données}

\begin{figure}[H]
    \centering
    \includegraphics[width=0.85\textwidth, keepaspectratio]{screen/6.png}
    \caption{Étape 1 -- Saisie des identifiants MariaDB (localhost, utilisateur dédié)}
\end{figure}

\subsubsection*{Étape 2 -- Test de connexion et sélection de la base}

\begin{figure}[H]
    \centering
    \includegraphics[width=0.85\textwidth, keepaspectratio]{screen/7.png}
    \caption{Étape 2 -- Connexion réussie, sélection de la base \texttt{glpi}}
\end{figure}

\subsubsection*{Étape 3 -- Initialisation de la base de données}

\begin{figure}[H]
    \centering
    \includegraphics[width=0.85\textwidth, keepaspectratio]{screen/8.png}
    \caption{Étape 3 -- Initialisation en cours (création des tables)}
\end{figure}

\begin{figure}[H]
    \centering
    \includegraphics[width=0.85\textwidth, keepaspectratio]{screen/8.1.png}
    \caption{Étape 3 -- Initialisation terminée à 100\,\%}
\end{figure}

\subsubsection*{Étapes 4 et 5 -- Télémétrie et informations complémentaires}

\begin{figure}[H]
    \centering
    \includegraphics[width=0.85\textwidth, keepaspectratio]{screen/9.png}
    \caption{Étape 4 -- Option de télémétrie (statistiques d'usage anonymes)}
\end{figure}

\begin{figure}[H]
    \centering
    \includegraphics[width=0.85\textwidth, keepaspectratio]{screen/10.png}
    \caption{Étape 5 -- Information sur le support GLPI-Network}
\end{figure}

\subsubsection*{Étape 6 -- Fin de l'installation}

\begin{figure}[H]
    \centering
    \includegraphics[width=0.85\textwidth, keepaspectratio]{screen/11.png}
    \caption{Étape 6 -- Installation terminée, identifiants par défaut affichés}
\end{figure}

\subsubsection*{Première connexion}

\begin{figure}[H]
    \centering
    \includegraphics[width=0.85\textwidth, keepaspectratio]{screen/12.png}
    \caption{Page de connexion GLPI après installation}
\end{figure}

\subsection{Réponses aux questions}

\begin{description}
    \item[Q1 -- Pourquoi créer un utilisateur MariaDB dédié plutôt qu'utiliser \texttt{root} ?]
    Le compte \texttt{root} possède des droits complets sur toutes les bases de données. En cas de compromission de GLPI, un attaquant aurait accès à tout le serveur MariaDB. Un utilisateur dédié avec des droits limités à la seule base \texttt{glpi} réduit considérablement la surface d'attaque.

    \item[Q2 -- À quoi sert l'encodage \texttt{utf8mb4} ?]
    \texttt{utf8mb4} est l'encodage UTF-8 complet sur 4 octets. Il permet de stocker tous les caractères Unicode, y compris les émojis et les caractères asiatiques, sans risque de troncature ou d'erreur. Il est recommandé pour toute application moderne.

    \item[Q3 -- Différence entre installation TP local et installation de production ?]
    En local (TP), on utilise des mots de passe simples, pas de HTTPS, et les données ne sont pas critiques. En production, on sécurise les accès (HTTPS, mots de passe forts, pare-feu), on met en place des sauvegardes régulières, des mises à jour automatiques et une supervision, car des données réelles d'entreprise sont en jeu.
\end{description}

\newpage

% ============================================================
\section{Exercice 3 -- Création des utilisateurs, rôles et catégories}
% ============================================================

\subsection{Utilisateurs créés}

J'ai créé les trois comptes suivants dans GLPI (\textit{Administration > Utilisateurs}) :

\begin{center}
\begin{tabularx}{0.95\textwidth}{l l X}
\toprule
\textbf{Nom du compte} & \textbf{Rôle} & \textbf{Utilisation dans le TP} \\
\midrule
\texttt{enseignant1} & Utilisateur (Self-Service) & Signale des incidents ou demandes \\
\texttt{admin1} & Utilisateur (Self-Service) & Signale une demande administrative \\
\texttt{technicien1} & Technicien & Traite les tickets \\
\bottomrule
\end{tabularx}
\end{center}

\subsection{Catégories de tickets créées}

Les catégories suivantes ont été créées dans \textit{Configuration > Intitulés > Catégories de tickets} :

\begin{itemize}
    \item Compte et accès ;
    \item Réseau et Wi-Fi ;
    \item Matériel ;
    \item Logiciel ;
    \item Messagerie.
\end{itemize}

\subsection{Organisation et intérêt}

Les utilisateurs (enseignant1, admin1) disposent du profil \textbf{Self-Service} : ils peuvent créer des tickets et suivre l'avancement des leurs, sans accès à l'administration. Le compte \texttt{technicien1} dispose du profil \textbf{Technicien} : il consulte, prend en charge et résout les tickets qui lui sont assignés. Cette séparation des droits garantit la sécurité et la cohérence du système.

\subsection{Réponses aux questions}

\begin{description}
    \item[Q1 -- Différence entre requérant et technicien ?]
    Le \textbf{requérant} est l'utilisateur qui signale un problème et suit son ticket. Il n'a pas de visibilité sur les autres tickets ni sur les actions techniques. Le \textbf{technicien} prend en charge les tickets, effectue les interventions et rédige les solutions. Il a accès à l'ensemble de la file d'attente.

    \item[Q2 -- Pourquoi catégoriser les tickets ?]
    La catégorisation permet de trier rapidement les demandes, d'identifier les problèmes récurrents, de router les tickets vers le bon technicien compétent, et de produire des statistiques par type d'incident pour anticiper les besoins.

    \item[Q3 -- À quoi servent les profils et les droits dans GLPI ?]
    Les profils définissent précisément ce que chaque type d'utilisateur peut voir et faire dans GLPI. Cela évite qu'un utilisateur lambda accède à des données confidentielles, modifie des configurations ou consulte les tickets des autres, tout en permettant au technicien et à l'administrateur de travailler efficacement.
\end{description}

\newpage

% ============================================================
\section{Exercice 4 -- Création et traitement des tickets}
% ============================================================

\subsection{Tickets créés}

J'ai créé les quatre tickets suivants, conformément aux scénarios du TP :

\begin{center}
\small
\begin{tabularx}{\textwidth}{c X l l l X}
\toprule
\textbf{N°} & \textbf{Titre} & \textbf{Type} & \textbf{Catégorie} & \textbf{Priorité} & \textbf{Solution résumée} \\
\midrule
1 & Impossible d'ouvrir la session Windows & Incident & Compte et accès & Haute & Réinitialisation du mot de passe utilisateur via l'AD \\
2 & Accès imprimante salle 203 pour M. Dupont & Demande & Matériel & Moyenne & Ajout du droit d'impression dans les paramètres réseau \\
3 & Wi-Fi indisponible en salle B12 & Incident & Réseau et Wi-Fi & Critique & Redémarrage du point d'accès défaillant AP-B12 \\
4 & Installation de Visual Studio Code & Demande & Logiciel & Basse & Installation réalisée après validation du responsable \\
\bottomrule
\end{tabularx}
\end{center}

\subsection{Détail des tickets}

\subsubsection*{Ticket 1 -- Impossible d'ouvrir la session Windows}
\begin{description}
    \item[Description :] L'utilisateur \texttt{enseignant1} signale qu'il ne peut plus se connecter à son poste Windows depuis ce matin. Le message affiché est « Le nom d'utilisateur ou le mot de passe est incorrect ».
    \item[Suivi technique :] Vérification que le verrou Majuscules n'est pas activé, puis test du compte sur un autre poste confirmant le blocage. Réinitialisation du mot de passe via l'administrateur.
    \item[Solution :] Mot de passe réinitialisé. L'utilisateur a pu se connecter. Il lui est conseillé de choisir un mot de passe mémorisable mais robuste.
\end{description}

\subsubsection*{Ticket 2 -- Accès imprimante salle 203}
\begin{description}
    \item[Description :] M. Dupont, nouvel enseignant, demande l'accès à l'imprimante réseau de la salle 203 (\texttt{IMP-203}) pour imprimer ses cours.
    \item[Suivi technique :] Ajout du compte dans le groupe de sécurité « Impression-Salle203 » sur le serveur.
    \item[Solution :] L'accès à l'imprimante a été accordé. L'utilisateur a été informé de la procédure d'installation du pilote sur son poste.
\end{description}

\subsubsection*{Ticket 3 -- Wi-Fi indisponible en salle B12}
\begin{description}
    \item[Description :] Le Wi-Fi est totalement indisponible dans la salle B12. Plusieurs étudiants et l'enseignant ne peuvent pas accéder à Internet depuis leurs appareils.
    \item[Suivi technique :] Connexion à la console de supervision réseau : le point d'accès AP-B12 est signalé hors ligne. Redémarrage à distance depuis le contrôleur Wi-Fi. Test de reconnexion concluant.
    \item[Solution :] Le point d'accès AP-B12 a été redémarré à distance. Le Wi-Fi est rétabli dans la salle. À surveiller (possible défaillance matérielle à venir).
\end{description}

\subsubsection*{Ticket 4 -- Installation de Visual Studio Code}
\begin{description}
    \item[Description :] Un étudiant demande l'installation de Visual Studio Code sur son poste pour les cours de développement.
    \item[Suivi technique :] Vérification de la validation du logiciel dans la liste des logiciels autorisés. Téléchargement depuis le site officiel et installation silencieuse.
    \item[Solution :] Visual Studio Code installé et testé. L'étudiant a été informé de son installation.
\end{description}

\subsection{Réponses aux questions}

\begin{description}
    \item[Q1 -- Différence entre urgence, impact et priorité ?]
    L'\textbf{urgence} représente la vitesse à laquelle il faut intervenir selon le ressenti de l'utilisateur. L'\textbf{impact} mesure le nombre de personnes ou de services affectés par le problème. La \textbf{priorité} est calculée automatiquement par GLPI en combinant les deux : un incident urgent avec un fort impact sera traité en premier.

    \item[Q2 -- Pourquoi le Wi-Fi en salle est plus prioritaire qu'une installation logicielle ?]
    Le Wi-Fi indisponible en salle bloque simultanément l'accès à Internet pour un groupe entier d'étudiants et un enseignant, impactant directement le déroulement du cours. L'installation d'un logiciel ne concerne qu'un seul utilisateur et n'est pas bloquante immédiatement.

    \item[Q3 -- Pourquoi ajouter un suivi technique dans un ticket ?]
    Le suivi technique permet de tracer toutes les actions réalisées : ce qui a été vérifié, testé, modifié. Il est indispensable pour la communication entre techniciens sur un même ticket, et constitue un historique précieux pour diagnostiquer des incidents similaires à l'avenir.

    \item[Q4 -- Risque de clôturer sans solution détaillée ?]
    Sans solution documentée, l'équipe ne peut pas capitaliser sur l'expérience acquise. Si l'incident se reproduit, le technicien (ou un collègue) devra repartir de zéro. De plus, l'utilisateur ne comprend pas ce qui a été fait et ne peut pas prévenir une récidive.
\end{description}

\newpage

% ============================================================
\section{Exercice 5 -- Base de connaissances}
% ============================================================

\subsection{Article rédigé}

À partir du Ticket 1, j'ai rédigé l'article de base de connaissances suivant dans \textit{Outils > Base de connaissances}.

\vspace{0.5em}
\noindent\textbf{Titre de l'article :} \textit{Que faire lorsqu'un utilisateur ne peut plus ouvrir sa session Windows ?}

\vspace{0.5em}
\noindent\textbf{Problème rencontré :}\\
L'utilisateur ne peut pas se connecter à son poste Windows. Le message affiché est « Nom d'utilisateur ou mot de passe incorrect ».

\vspace{0.5em}
\noindent\textbf{Symptômes observables :}
\begin{itemize}
    \item Impossible de saisir les identifiants sur l'écran de connexion.
    \item Message d'erreur au moment de la validation.
    \item Le problème persiste après redémarrage du poste.
\end{itemize}

\vspace{0.5em}
\noindent\textbf{Procédure de résolution :}
\begin{enumerate}
    \item Vérifier que la touche Verr. Maj n'est pas activée.
    \item Vérifier que la disposition du clavier est correcte (FR, pas EN).
    \item Tester la connexion sur un autre poste pour isoler le problème.
    \item Si l'échec persiste, contacter le technicien pour réinitialiser le mot de passe.
    \item Après réinitialisation, demander à l'utilisateur de changer son mot de passe à la première connexion.
    \item Documenter l'intervention dans GLPI et clôturer le ticket.
\end{enumerate}

\vspace{0.5em}
\noindent\textbf{Vérifications finales :}\\
S'assurer que l'utilisateur peut se connecter normalement et accéder à ses fichiers sans erreur.

\vspace{0.5em}
\noindent\textbf{Remarque de prévention :}\\
Inviter les utilisateurs à noter leurs identifiants dans un gestionnaire de mots de passe sécurisé pour éviter ce type d'incident.

\subsection{Réponses aux questions}

\begin{description}
    \item[Q1 -- Comment la base de connaissances réduit-elle le nombre de tickets ?]
    Si un utilisateur rencontre un problème déjà documenté, il peut consulter lui-même la procédure et résoudre son incident sans contacter le support. Cela libère du temps aux techniciens et autonomise les utilisateurs.

    \item[Q2 -- Différence entre suivi technique interne et article visible par les utilisateurs ?]
    Le suivi technique est réservé aux techniciens : il peut contenir des informations sensibles (identifiants, configurations réseau, commandes système) qui ne doivent pas être divulguées. L'article de la base de connaissances est rédigé dans un langage accessible aux utilisateurs finaux, sans informations sensibles.

    \item[Q3 -- Pourquoi rédiger une procédure simple et compréhensible ?]
    La base de connaissances est destinée à des utilisateurs non techniciens. Si la procédure est trop complexe ou utilise du jargon, l'utilisateur ne pourra pas l'appliquer seul. Une procédure claire, avec des étapes numérotées et un vocabulaire adapté, garantit son efficacité.
\end{description}

\newpage

% ============================================================
\section{Exercice 6 -- Analyse du fonctionnement du helpdesk}
% ============================================================

La mise en place d'un outil de ticketing comme GLPI représente un changement structurant pour une organisation. Par rapport à une gestion par mail ou par téléphone, GLPI apporte une centralisation complète des demandes : chaque incident ou demande est formalisé, traçable et consultable à tout moment par les personnes concernées. Plus aucune demande ne peut être « oubliée dans une boîte mail » ou perdue dans un couloir.

Pour qu'un ticket soit réellement exploitable, il doit contenir a minima un titre explicite, une description précise du problème, un type (incident ou demande), une catégorie, une urgence, un impact, une priorité et le nom du requérant. Sans ces informations, le technicien ne peut pas agir efficacement.

Une mauvaise catégorisation des tickets entraîne des problèmes concrets : difficulté à assigner les tickets aux bons techniciens, statistiques faussées, et impossibilité de détecter les incidents récurrents par domaine. Si les tickets sont mal triés, l'équipe ne peut pas anticiper les besoins ni améliorer les procédures.

L'historique des tickets est une ressource précieuse. Il permet d'identifier les pannes récurrentes, de mesurer les délais de traitement, d'évaluer la charge de travail de chaque technicien et de justifier des demandes d'investissement matériel ou logiciel auprès de la direction.

Enfin, pour utiliser GLPI dans un vrai établissement, plusieurs précautions sont indispensables : activer le chiffrement HTTPS, appliquer une politique de mots de passe robuste, mettre en place des sauvegardes régulières de la base de données, gérer les droits d'accès avec rigueur, et sensibiliser les utilisateurs à l'outil pour en garantir l'adoption.

\newpage

% ============================================================
\section{Conclusion}
% ============================================================

Ce TP m'a permis de découvrir et de pratiquer l'ensemble du cycle de vie d'un ticket dans GLPI, depuis l'installation de l'outil jusqu'à la capitalisation des solutions. Les principales compétences acquises sont :

\begin{itemize}
    \item installer une instance GLPI locale sous WAMP ;
    \item créer des utilisateurs avec des profils et des droits distincts ;
    \item structurer un service support avec des catégories de tickets ;
    \item créer, traiter et clôturer des tickets d'incidents et de demandes ;
    \item documenter les interventions et alimenter une base de connaissances ;
    \item analyser l'apport professionnel d'un outil de ticketing.
\end{itemize}

\textbf{Difficultés rencontrées :} \textit{(À compléter en fin de TP selon les problèmes rencontrés.)}

\end{document}
